\section{Methodology}
\label{sec:methodology}

All numbers come from one rig: a Dell T5820 (Xeon W-2140B) holding two
Quadro RTX 6000 GPUs, both at PCIe 3.0~x16, bridged by a two-link NV2
NVLink, driver 610.43.02, CUDA toolkit 13.3.33, ECC disabled on both
devices, CPU governor pinned to \texttt{performance}. The table is a
snapshot of this toolchain and driver pairing, not a living document. The
append-only results layout admits later datasets without disturbing this
one.

\subsection{Clock contract}
\label{sec:clocks}

Cycle-domain rows require a fixed SM clock. The harness prechecks a
1455~MHz lock (\texttt{nvidia-smi -lgc}) at startup and exits otherwise.
The memory clock cannot be locked on this part: \texttt{-lmc} reports
unsupported and the applications-clock interface is deprecated in this
driver. What holds instead, measured on both GPUs, is that CUDA compute
always executes in the P2 performance state with the memory clock at
6500~MHz. The harness ramps into P2 with a warmup kernel, verifies the
memory clock under load, and samples both clocks around every repetition,
rejecting any where either is off target. Bandwidth expectations are
therefore stated against the P2 memory clock: a 624~GB/s DRAM ceiling
rather than the 672 implied by the unreachable P0 point.

\subsection{Timing}

Latency rows use \texttt{clock64()} around dependent chains, with loop
overhead cancelled by a doubled-trip-count slope and the loop body
verified by size to fit the L0 instruction cache. Pointer-chase rows
embed full 64-bit pointers (or byte offsets, in shared and constant
windows) so the timed step is a single dependent load with no address
arithmetic \citep{wong2010demystifying,mei2017dissecting}. Throughput
rows are also timed in actual SM cycles, by one block on one SM:
wall-clock event timing normalised by the nominal locked clock showed a
0.3--1.0\% between-run spread that grew with region length (real-clock
thermal sag under the lock), where on-device cycle counts are immune.
Host-domain rows (launches, events, NCCL calls) use
\texttt{steady\_clock} over a thousand repetitions with a device-event
cross-check, reported as median with the 10th and 90th percentiles.

\subsection{SASS verification}
\label{sec:sass}

Operation selection is pinned by PTX inline assembly, and
\texttt{check\_sass.py} disassembles every benchmark binary, locates each
timed loop, and enforces that it contains the intended instructions and
nothing unexplained. The gate is not a formality. Working through the
table, \texttt{ptxas} or the NVVM front end variously: constant-folded
integer chains built on literal operands; moved warp-uniform integer
chains onto the uniform datapath (\texttt{UPOPC} on \texttt{UR}
registers); strength-reduced dependent add chains into \texttt{IMAD} and
\texttt{LEA} through both a dead carry-out pin and cross-coupled
accumulators; deleted six of eight independent accumulator chains whose
results a sink did not consume; folded \texttt{rcp.approx} compositions
as exact algebra, twice; contracted a deliberately separate add into an
\texttt{FFMA}; lowered an f16 widening conversion to \texttt{HADD2}
rather than the expected \texttt{F2F}; and expanded a 64-bit remainder
into a division subroutine call whose float-reciprocal emulation put
\texttt{MUFU} and conversion traffic inside a bandwidth loop. Every one
of these was caught by counting emitted mnemonics against the design,
not by a timing looking wrong. Several produced plausible numbers for
the wrong instruction stream. Both gates carry negative tests: a
deliberately unlocked clock and a deliberately miscompiled kernel must
fail.

\subsection{Statistics and gate taxonomy}

Each row is the median of ten repetitions, with the coefficient of
variation published. Every published row further requires at least two
independent process invocations on each GPU, and the between-run spread
must not exceed the within-run variation beyond a domain floor (0.1\% for
cycle rows, 0.5\% for wall-clock bandwidths, 5\% for host-domain times).
SM-domain rows must agree across the two physical GPUs within their
combined variation. Interconnect rows differ by design and carry
direction in the variant. Verification gates are of two kinds, and
confusing them is how benchmark suites quietly tune themselves to the
literature: \emph{methodology-sanity} gates (two independent methods must
agree) block publication on failure, while \emph{prior-comparison} gates
(against \citealp{jia2019turing} or whitepaper geometry) publish the
deviation whatever it is. The table exists to report true deviations.

\subsection{Pre-registration}
\label{sec:prereg}

The coverage manifest and measurement policy were committed to the public
repository before any benchmark ran, and the three decision-grade
hypotheses of \cref{sec:hypotheses}, with their predictions,
operationalisations, and decision rules, were committed before their
data existed (repository commit \texttt{a4fdd73}, 2026-06-10, preceding
every measurement it governs). The commit history is the receipt. One
registered rule fired mid-study and was honoured: a composed-prediction
gate failed at its registered tolerance, and the affected comparison was
blocked rather than re-derived, until independently measured constituents
allowed a documented remediation (\cref{sec:hypotheses-outcomes}).

\subsection{A worked row}

The discipline end to end, for \texttt{alu.ffma.lat}. A 128-deep
dependent \texttt{fma.rn.f32} chain on runtime operands compiles to
exactly 128 \texttt{FFMA} in the timed loop, and the gate passes. Ten
repetitions at two trip counts give a slope of 4.07 cycles with zero
within-run variation, and four invocations across both GPUs agree within
the 0.1\% floor. Against the 4-cycle prior
\citep[tab.~4.1]{jia2019turing} the row publishes as 4.07, deviation
$+1.8\%$, flag \texttt{ok}, with provenance naming the benchmark file,
commit, and results lines.
